\documentclass[a4paper]{article}
\usepackage{cmap}  
\usepackage{vntex}
%\usepackage[english,vietnam]{babel}
%\usepackage[utf8]{inputenc}
%\usepackage[utf8]{inputenc}
%\usepackage[francais]{babel}
\usepackage{a4wide,amssymb,epsfig,latexsym,multicol,array,hhline,fancyhdr}

\usepackage{amsmath}
\usepackage{lastpage}
%\usepackage[lined,boxed,commentsnumbered]{algorithm2e}
\usepackage{enumerate}
\usepackage{color}
\usepackage{graphicx}							% Standard graphics package
\usepackage{array}
\usepackage{tabularx, caption}
\usepackage{multirow}
\usepackage{multicol}
\usepackage{rotating}
\usepackage{graphics}
\usepackage{geometry}
\usepackage{setspace}
\usepackage{epsfig}
\usepackage{tikz}
\usetikzlibrary{arrows,snakes,backgrounds}
\usepackage[unicode]{hyperref}
\hypersetup{urlcolor=blue,linkcolor=black,citecolor=black,colorlinks=true} 
%\usepackage{pstcol} 								% PSTricks with the standard color package


\usepackage{pifont}% http://ctan.org/pkg/pifont
\newcommand{\cmark}{\ding{51}}%
\newcommand{\xmark}{\ding{55}}%

%\usepackage{fancyhdr}
\setlength{\headheight}{40pt}
\pagestyle{fancy}
\fancyhead{} % clear all header fields
\fancyhead[L]{
 \begin{tabular}{rl}
    \begin{picture}(25,15)(0,0)
    \put(0,-8){\includegraphics[width=8mm, height=8mm]{LogoBK.jpg}}
    %\put(0,-8){\epsfig{width=10mm,figure=hcmut.eps}}
   \end{picture}&
	\begin{tabular}{l}
		\textbf{\bf \ttfamily Faculty of Computer Science and Engineering}\\
		\textbf{\bf \ttfamily Bach Khoa University}
	\end{tabular} 	
 \end{tabular}
}
\fancyhead[R]{
	\begin{tabular}{l}
		\tiny \bf \\
		\tiny \bf 
	\end{tabular}  }
\fancyfoot{} % clear all footer fields
\fancyfoot[L]{\scriptsize \ttfamily Report Cryptography and Network Security - Year 2017-2018}
\fancyfoot[R]{\scriptsize \ttfamily Page {\thepage}/\pageref{LastPage}}
\renewcommand{\headrulewidth}{0.3pt}
\renewcommand{\footrulewidth}{0.3pt}


%%%
%\setcounter{secnumdepth}{4}
%\setcounter{tocdepth}{3}
%\makeatletter
%\newcounter {subsubsubsection}[subsubsection]
%\renewcommand\thesubsubsubsection{\thesubsubsection .\@alph\c@subsubsubsection}
%\newcommand\subsubsubsection{\@startsection{subsubsubsection}{4}{\z@}%
%                                     {-3.25ex\@plus -1ex \@minus -.2ex}%
%                                     {1.5ex \@plus .2ex}%
%                                     {\normalfont\normalsize\bfseries}}
%\newcommand*\l@subsubsubsection{\@dottedtocline{3}{10.0em}{4.1em}}
%\newcommand*{\subsubsubsectionmark}[1]{}
%\makeatother



\begin{document}

\begin{titlepage}
\begin{center}
Faculty of Computer Science and Engineering\\
Ho Chi Minh City University of Technology
\end{center}

\vspace{.51cm}

\begin{figure}[h!]
\begin{center}
\includegraphics[width=3cm]{LogoBK.jpg}
\end{center}
\end{figure}

\vspace{0.5cm}


\begin{center}
\begin{tabular}{c}
\multicolumn{1}{l}{\textbf{{\Large \textcolor{blue}{Xử Lý Ảnh Và Thị Giác Máy Tính}}}}\\
~~\\
\hline
\\
\multicolumn{1}{l}{\textbf{{\Large \textcolor{blue}{Đề Tài}}}}\\
\\
\textbf{{\Huge \textcolor{blue}{Nhận Diện Con Người RealTime}}}\\
\\
\hline
\end{tabular}
\end{center}

\vspace{3cm}

\begin{table}[h]
\begin{tabular}{rrll}
\hspace{5 cm} & GV Hướng Dẫn: & Trần Tuấn Anh &\\
& Author: & Nguyễn Thành Đạt & 1510700 \\
& & ? & ? \\ 
& & ?& ?\\ 
\end{tabular}
\end{table}

\begin{center}
{\footnotesize Bach Khoa, 9/2018}
\end{center}
\end{titlepage}


%\thispagestyle{empty}
\newpage
\tableofcontents
\newpage

%%%%%%%%%%%%%%%%%%%%%%%%%%%%%%%%%
\section{Mô tả đề tài}
\hspace{6mm} Đề tài nhóm hướng tới là nhận dạng đối tượng là con người qua camera tại thời gian thực thông qua một số đặc trưng của hình dáng (body) lẫn khuôn mặt.

\section{Các Bước Thực Hiện}
\subsection{Chuẩn bị dữ liệu}
\hspace{6mm} Nhóm sẽ xác định và sử dụng tập data có sẵn.
\subsection{•}
\hspace{6mm}Sau đây chúng tôi sẽ giới thiệu cho các bạn 3 phương pháp mã hóa thông dụng cũng như nguyên lý mã hóa của chúng:

\section{Kết luận }


\section{Reference}
%%%%%%%%%%%%%%%%%%%%%%%%
%%%%%%%%%%%%%%%%%%%%%%%%%%%%%%%%%


 
\end{document}

